% 第5节：几何产生的规范理论
\section{来自几何的规范理论}
\label{sec:gauge}

\subsection{规范对称性的几何起源}

在标准模型中，规范群 $U(1) \times SU(2) \times SU(3)$ 是基于实验观察假设的。我们的理论通过扭曲覆盖映射的核分解为这些对称性提供了几何起源。

回顾第 \ref{sec:math} 节：
\begin{equation}
    \ker(\rho^\tau) \cong Z_2^5
\end{equation}

这个五重 $Z_2$ 结构通过以下机制分解为规范群。

\subsection{核分解}

五个独立的 $Z_2$ 因子可以分组为：
\begin{itemize}
    \item 1个因子 $\to$ $U(1)$：超荷
    \item 2个因子 $\to$ $SU(2)$：弱同位旋  
    \item 3个因子 $\to$ $SU(3)$：色（强相互作用）
\end{itemize}

分解保持乘积结构：
\begin{equation}
    Z_2 \times Z_2^2 \times Z_2^3 \cong Z_2^5 \to U(1) \times SU(2) \times SU(3)
\end{equation}

\subsection{电磁学：U(1)}

电磁规范群 $U(1)_{em}$ 从单个 $Z_2$ 因子中涌现。规范势为：
\begin{equation}
    A_\mu = \tau_0 \cdot \text{Tr}(\gamma_5 \gamma_\mu \Sigma)
\end{equation}
其中 $\Sigma$ 是带挠率的自旋联络。

场强张量为：
\begin{equation}
    F_{\mu\nu} = \partial_\mu A_\nu - \partial_\nu A_\mu + \mathcal{O}(\tau_0^2)
\end{equation}

在 $\tau_0$ 的领先阶，这复现麦克斯韦方程组：
\begin{align}
    \partial_\mu F^{\mu\nu} &= J^\nu \\
    \partial_{[\mu} F_{\nu\rho]} &= 0
\end{align}

\subsection{弱相互作用：SU(2)}

弱相互作用规范群 $SU(2)_L$ 从两个 $Z_2$ 因子中涌现。规范势为：
\begin{equation}
    W^a_\mu = \tau_0 \cdot \text{Tr}(\sigma^a \gamma_\mu \Sigma)
\end{equation}
其中 $\sigma^a$ 是作用在弱同位旋空间上的泡利矩阵。

场强张量为：
\begin{equation}
    W^a_{\mu\nu} = \partial_\mu W^a_\nu - \partial_\nu W^a_\mu + g \epsilon^{abc} W^b_\mu W^c_\nu
\end{equation}

耦合常数 $g$ 几何地确定：
\begin{equation}
    g = \tau_0 \cdot \sqrt{2}
\end{equation}

对于 $\tau_0 = 10^{-6}$：
\begin{equation}
    g \approx 1.4 \times 10^{-6}
\end{equation}

这是裸耦合；重整化群跑动给出测量值 $g(m_Z) \approx 0.65$。

\subsection{强相互作用：SU(3)}

强相互作用规范群 $SU(3)_C$ 从三个 $Z_2$ 因子中涌现。规范势为：
\begin{equation}
    G^\alpha_\mu = \tau_0 \cdot \text{Tr}(\lambda^\alpha \gamma_\mu \Sigma)
\end{equation}
其中 $\lambda^\alpha$ 是盖尔曼矩阵。

场强张量为：
\begin{equation}
    G^\alpha_{\mu\nu} = \partial_\mu G^\alpha_\nu - \partial_\nu G^\alpha_\mu + g_s f^{\alpha\beta\gamma} G^\beta_\mu G^\gamma_\nu
\end{equation}

强耦合为：
\begin{equation}
    g_s = \tau_0 \cdot \sqrt{3}
\end{equation}

\subsection{CKM和PMNS矩阵}

夸克混合（CKM）和轻子混合（PMNS）矩阵源于内部空间中规范群生成元的相对方向。

对于CKM矩阵：
\begin{equation}
    V_{CKM} = \exp\left(i \theta_{12} \sigma_2 \otimes \sigma_2 + i \theta_{13} \sigma_1 \otimes \sigma_3 + i \theta_{23} \sigma_3 \otimes \sigma_1\right)
\end{equation}

混合角由以下确定：
\begin{align}
    \theta_{12} &\approx \arctan\left(\frac{m_s}{m_d}\right) \cdot \tau_0^{1/2} \\
    \theta_{13} &\approx \frac{m_u}{m_t} \cdot \tau_0 \\
    \theta_{23} &\approx \arctan\left(\frac{m_c}{m_s}\right) \cdot \tau_0^{1/2}
\end{align}

这些给出了CKM角的正确数量级。

\subsection{费米子质量}

费米子质量源于与希格斯样场（在我们的理论中是几何的）的耦合：
\begin{equation}
    m_f = y_f \cdot v \cdot \tau_0^{n_f}
\end{equation}

其中 $n_f$ 依赖于费米子代：
\begin{itemize}
    \item 第1代：$n_f = 1$
    \item 第2代：$n_f = 2$  
    \item 第3代：$n_f = 3$
\end{itemize}

这自然地解释了代之间的质量层次。

\subsection{耦合常数统一}

规范耦合随能量尺度跑动并在GUT尺度 $M_{GUT} \sim 10^{16}$ GeV 处统一：
\begin{equation}
    g_1(M_{GUT}) = g_2(M_{GUT}) = g_3(M_{GUT}) = \sqrt{4\pi\tau_0}
\end{equation}

这是因为在高能时所有规范对称性都源自相同的几何结构。

\subsection{小结}

我们框架中的规范理论：
\begin{itemize}
    \item $U(1) \times SU(2) \times SU(3)$ 从 $Z_2^5$ 核分解中涌现
    \item 耦合常数由 $\tau_0$ 确定：$g_i \propto \tau_0^{n_i/2}$
    \item CKM/PMNS混合角从质量比预言
    \item 费米子质量层次由代依赖的 $\tau_0$ 幂次解释
    \item 规范耦合在GUT尺度处统一
\end{itemize}

这完成了所有基本力的几何统一。
